\chapter{Sources, Evidence, and Research Corpus}
\label{chap:sources-evidence}

\classification{Evidence status: Corpus-description chapter. Confidence baseline: High for repository counts and catalogue structure; variable for source completeness and authenticity. This chapter describes the evidence base and its limits. It does not interpret the evidence or evaluate case hypotheses.}

\section{Overview of the Research Corpus}

This chapter describes the documentary and evidentiary foundation of the investigation. It is concerned with what the project holds, what it has catalogued, how the material is classified, and where the documentary record remains incomplete. Factual reconstruction begins in later chapters.

The frozen research corpus contains 24 entries in the source register, 23 entries in the primary-document inventory, 63 audited bibliography records, 18 evidence-register items, 9 image-catalogue leads, and 10 planned figure entries requiring rights or replacement review. These counts describe the repository as prepared for manuscript drafting, not the total universe of possible Somerton Man records.

The corpus combines official record references, public reproductions of primary-source material, contemporary newspaper reporting, modern reporting, academic and professional literature, institutional guides, image leads, physical-evidence catalogue entries, and project governance records. The strongest case-specific documentary anchors are the 1949 coronial inquest record reference, the 1958 later inquest record reference, the acquired public reproductions of those inquest files, police-file reproductions, death-certificate material, fingerprint and dental charts, and the preserved image of the Rubaiyat code markings.

The repository also records sources that were not acquired directly from original custodians. Most acquired primary-source files are public reproductions from Derek Abbott's primary-source archive rather than certified copies from the original repository. The archive is therefore usable for controlled analysis, but direct official copies remain preferable where available.

\begin{table}[htbp]
\centering
\caption{Research corpus summary at manuscript drafting freeze}
\label{tab:sources-corpus-summary}
\begin{tabularx}{\textwidth}{p{0.34\textwidth}rX}
\toprule
Repository asset & Count & Evidentiary role \\
\midrule
Source-register entries & 24 & Bibliographic, reliability, and source-priority control records. \\
Primary-document inventory records & 23 & Acquired or located primary-source reproductions and source leads. \\
Audited bibliography records & 63 & End-reference records generated from source, document, and Mission 11 registers. \\
Evidence-register items & 18 & Physical, documentary, and forensic evidence items or leads. \\
Image-catalogue leads & 9 & Candidate visual material requiring source and rights review. \\
Planned figure entries & 10 & Publication-planning records; no rights-restricted image is inserted in this chapter. \\
\bottomrule
\end{tabularx}
\end{table}

\section{Primary Sources}

Primary-source material in this project includes coronial records, police records, death-certificate material, contemporary official statements, evidence images, fingerprint and dental charts, directory extracts, and public legal or government records. These materials carry the highest evidentiary value when their authenticity, custody, and completeness can be established.

The primary-document inventory contains 23 records. Nineteen are recorded as acquired local digital copies. Four are recorded as located but not acquired locally during the primary-source acquisition pass. The inventory classifies 22 records as later reproductions and one record as a certified copy. That authenticity profile is central to later source weighting.

The highest-priority case documents are the 1949 inquest reproduction, the 1958 later inquest reproduction, the unknown man's death-certificate copy, Criminal Investigation Branch reports, a 1949 police-file extract, fingerprint and tooth charts, and the Rubaiyat code image. These records provide a working primary-source layer, but several remain reproductions whose original custodian copies should be obtained or compared where possible.

\begin{table}[htbp]
\centering
\caption{Primary-document inventory by acquisition and authenticity status}
\label{tab:sources-primary-document-status}
\begin{tabularx}{\textwidth}{p{0.36\textwidth}rX}
\toprule
Status field & Count & Meaning for use \\
\midrule
Acquired local digital copy & 19 & Available in the project archive for review, transcription, indexing, and source mapping. \\
Located but not acquired & 4 & Identified at source but not preserved locally in the acquisition pass. \\
Later reproduction & 22 & Usable as a source reproduction, but not equivalent to a direct official custodian copy. \\
Certified copy & 1 & Higher authenticity status within the acquired inventory. \\
\bottomrule
\end{tabularx}
\end{table}

The 1949 and 1958 inquest reproductions are especially important because they gather coronial, medical, police, witness, and administrative material in a form that can be indexed at page level. The repository treats those reproductions as high-priority working sources, while preserving the need for official archive copies and page-for-page comparison.

Police-source material is present but incomplete. The project has acquired or catalogued Criminal Investigation Branch reports, a police-file extract, police-gazette material, and a police/coroner report in related material. The current SAPOL open investigation file, original property registers, and original police evidence inventories remain unavailable or restricted in the research corpus.

Medical and forensic primary material is present through inquest testimony, forensic identifiers, the fingerprint chart, the tooth chart, and medical registers generated from primary-source review. Complete original laboratory records, retained-sample documentation, and post-2021 forensic reports were not acquired.

\section{Secondary Sources}

Secondary sources are used to provide technical context, source discovery, historical interpretation, and leads for primary-source retrieval. They do not replace original records when an original or official record exists.

The source register includes academic and professional material, book-length research, institutional history material, modern reporting, and supplementary research tools. These sources are useful when they disclose provenance, identify archive pathways, summarize public developments, or provide disciplinary context. Their evidentiary value is lower when they mix interpretation with fact or repeat older claims without direct source support.

Modern identity, DNA, and genealogy reporting is treated as an update layer. It records what researchers, laboratories, news organizations, or public agencies have stated, but it does not convert researcher-led identification claims into an official identity determination unless an official coroner, SAPOL, or Forensic Science SA record is obtained and catalogued.

The prior analytical dossier is retained only as an internal prior-assessment source. It is useful for extracting claims, hypotheses, and research questions. It is not used as primary evidence for case facts.

\section{Archival Collections}

The repository records material or source pathways connected with State Records of South Australia, South Australia Police, Forensic Science SA, the National Archives of Australia, the National Library of Australia's Trove newspaper database, the National Film and Sound Archive of Australia, University of Adelaide-hosted source material, ABC reporting, Smithsonian reporting, and other publication or library pathways.

State Records of South Australia is the principal identified custodian pathway for the coronial inquest records. The 1949 file is identified as GRG 1/27 File 71/1949. The 1958 later file is identified as GRG 1/27 File 53/1958 and is recorded as requiring access through the Coroner's Court pathway.

Trove and the National Library of Australia provide contemporary newspaper records and public legal notices. These sources are valuable for contemporaneous reporting, public chronology, and source leads, but they remain below official transcripts and original case records in the source hierarchy.

The National Archives of Australia provides government and institutional context rather than direct case evidence in the current corpus. Its role is to support historical and institutional background where relevant, not to establish direct case linkage.

The University of Adelaide-hosted Abbott archive is a major practical source pathway for public reproductions. It supplies the acquired copies of the inquest reproductions and several related source files. Its value is high for locating and preserving material, but its files are treated as later reproductions unless independently confirmed by the original custodian.

\begin{table}[htbp]
\centering
\caption{Repository and archive roles in the corpus}
\label{tab:sources-archive-summary}
\begin{tabularx}{\textwidth}{p{0.30\textwidth}X p{0.24\textwidth}}
\toprule
Repository or pathway & Role in the corpus & Principal limitation \\
\midrule
State Records of South Australia & Identified custodian pathway for coronial files and archival case references. & Direct official copies not fully acquired in this project stage. \\
South Australia Police / Forensic Science SA & Official pathway for current investigation status, property records, exhumation, and forensic releases. & Current files and full forensic reports remain unavailable or restricted. \\
National Library of Australia / Trove & Contemporary newspaper reporting and public notices. & OCR and press-reporting limitations; lower weight than official records. \\
National Archives of Australia & Government and institutional context. & Contextual rather than direct case evidence in the present corpus. \\
University of Adelaide-hosted Abbott archive & Public reproductions and source leads for inquest files, police extracts, charts, and images. & Mostly later reproductions; original custody must be checked separately. \\
ABC, Smithsonian, and other reporting outlets & Modern reporting, OSINT updates, and source leads. & Secondary reporting unless tied to an original document or official statement. \\
\bottomrule
\end{tabularx}
\end{table}

\section{Evidence Classification}

The project classifies evidence by the role it plays in the investigation. The classification does not determine truth by itself. It determines how a record should be handled, weighted, and cross-referenced.

Documentary evidence includes inquest files, police-file reproductions, reports, certificates, directories, legal notices, and source indexes. Physical evidence includes objects associated with the deceased, clothing, suitcase material, the Tamam Shud slip, the Rubaiyat copy, tickets, cigarettes, matches, forensic samples, and related exhibit leads. Medical evidence includes autopsy testimony, pathology observations, toxicology statements, time-of-death statements, anatomical identifiers, fingerprint records, dental records, and death-mask or bust material where documented.

Witness evidence includes sworn testimony, police statements, newspaper-reported statements, later recollections, and official correspondence. Photographic evidence includes images of the deceased, the code, forensic identifiers, scene material, and exhibit images. Genealogical evidence includes DNA, family-tree, civil-record, and biographical material connected with the Webb attribution and other identity leads. Contextual evidence includes historical, institutional, transport, forensic-science, intelligence, and social background. Analytical evidence includes project-generated verification matrices, confidence registers, evidence-weighting tables, and structured analytic products.

\begin{table}[htbp]
\centering
\caption{Evidence classification system}
\label{tab:sources-evidence-classification}
\begin{tabularx}{\textwidth}{p{0.23\textwidth}X p{0.28\textwidth}}
\toprule
Class & Examples in the repository & Use constraint \\
\midrule
Documentary & Inquest files, police extracts, certificates, directories, legal notices, source indexes. & Prefer original or official copies where available. \\
Physical & Clothing, suitcase material, paper slip, Rubaiyat copy, tickets, cigarettes, matches, samples. & Separate object existence from custody and interpretation. \\
Medical & Autopsy testimony, pathology observations, toxicology statements, identifiers, dental and fingerprint material. & Distinguish observation, test result, limitation, and conclusion. \\
Witness & Inquest testimony, police statements, reported statements, later recollections. & Preserve page references and reliability limits. \\
Photographic & Deceased images, exhibit images, code image, charts, scene or location images. & Do not publish unless source and rights status are resolved. \\
Genealogical & DNA, family, civil-record, and biographical identity material. & Separate researcher-led attribution from official identification status. \\
Contextual & Historical, institutional, forensic, transport, media, and intelligence background. & Context does not establish direct case linkage by itself. \\
Analytical & Verification matrices, confidence registers, ACH, evidence-weighting and gap registers. & Derived project analysis; not a substitute for source evidence. \\
\bottomrule
\end{tabularx}
\end{table}

\section{Source Reliability}

The project weights sources by proximity to the original event, official status, authenticity, custody, and transparency. Original official case records receive the highest weight. Certified or archival reproductions receive high weight when provenance is clear. Contemporary official statements and legal records receive high weight for what they explicitly state. Contemporary newspapers receive moderate weight and require caution because of OCR errors, compression, and press simplification.

Researcher accounts, laboratory participant accounts, later journalism, books, podcasts, and commentary are useful for leads, context, and public chronology, but they require separation from the underlying primary documents. A source may be reliable for reporting that a researcher made a claim while remaining insufficient to prove the underlying case fact.

Document authenticity is a key limitation in the current corpus. The primary-document inventory records 22 of 23 entries as later reproductions. This does not make them unusable. It means that exact wording, page sequence, missing enclosures, image quality, and official custody should be checked when the point matters.

OCR is treated as a search aid rather than a final transcription authority. Page images, manual transcription, and page-level indexing are required where exact wording affects a claim, witness statement, medical observation, or exhibit description.

Missing and restricted records are part of the reliability picture. The repository records unavailable or restricted official copies, police files, property registers, forensic reports, large directory scans, and the Wells Rubaiyat scan. Those gaps are documented so that absence of material is not mistaken for negative evidence.

\section{Repository Structure}

The research corpus is organized into source, primary-document, evidence, person, location, timeline, verification, chronology, witness, exhibit, victimology, forensic-medical, physical-evidence, intelligence, OSINT, image, reference, and publication-readiness layers. Each layer records a different kind of evidentiary relationship.

Stable identifiers preserve the audit trail. Source records use `SRC-` identifiers, acquired or located documents use `DOC-` identifiers, evidence items use `EV-` identifiers, people use `PER-` identifiers, locations use `LOC-` identifiers, and events use event identifiers. Later manuscript statements can therefore point back to the research layer without relying on narrative memory.

The reference audit integrates source-register records, primary-document records, and Mission 11 OSINT records into the end-reference system. At the publication-readiness audit, all 63 reference records had no missing required fields, and bibliography generation had been prepared for Biber validation.

The image and figure layers are deliberately separated from the prose. The image catalogue records 9 image leads, and the figure-rights register records 10 planned figure entries. The repository freeze report records that figure insertion is not ready for final publication until rights, permission, or replacement status is resolved.

\section{Limitations of the Documentary Record}

The documentary record remains incomplete. The project has not acquired direct official copies of every relevant record from original custodians. Public reproductions are available for key materials, but official custodian comparison remains a future task.

Several high-value records are unavailable or restricted in the present corpus. These include the direct official copy of the 1949 inquest file, direct access to the restricted 1958 later inquest file, the current SAPOL open investigation file, original police property and evidence registers, and full post-2021 Forensic Science SA or SAPOL forensic reports.

Physical-evidence documentation is also incomplete. Many evidence-register entries have unknown current custodian status, fragmentary chain-of-custody information, or unresolved survival status. The repository therefore distinguishes between a documented object, the later custody of that object, and any interpretation based on that object.

The image record is constrained by rights and provenance. Candidate images may be important as evidence leads, but they are not inserted into the manuscript until publication rights and source reliability are resolved.

Historical limits cannot be removed by later writing. Some records may have been lost, never created, retained outside public access, or preserved only through later reproduction. The manuscript therefore treats the corpus as a controlled and expandable evidence base rather than a complete official archive.
